\section{Deciding whether a composite is a layout}\label{sec:decide}

While a composite can be natively emitted as is, sometimes we can simplify the codegen
if that composite is actually itself a layout.

Fix a map $f : \mathrm{Coord}(S) \to \Z$ on a shape $S = (n_1,\dots,n_m)$.
We take $S$ flat, since a nested shape has the coordinates of the flat shape
of its leaves, regrouped. Write $c = (c_1,\dots,c_m)$ for a coordinate and
$\mathbf{0}$ for the origin. Every layout sends $\mathbf{0}$ to $0$, and
composition and swizzle preserve this, so $f(\mathbf{0}) = 0$ for every
composite. Maps take values in $\Z$ throughout this section: strides are
integers, and Section~\ref{sec:restrict} subtracts a constant.

\begin{definition}[Refinement]\label{def:refine}
A \emph{refinement} of a flat shape $S = (n_1,\dots,n_m)$ is a shape
$S' = (S_1,\dots,S_m)$ with each $S_i$ flat and $\size S_i = n_i$. A
coordinate of $S_i$ is identified with its flattening in $[0,n_i)$
(Definition~\ref{def:layout}), so $\mathrm{Coord}(S')$ and $\mathrm{Coord}(S)$
are the same set. A map $f : \mathrm{Coord}(S) \to \Z$ \emph{is the layout}
$S'{:}D'$ when
\[
f(c_1,\dots,c_m) \;=\;
(S'{:}D')\big(\mathrm{unflatten}_{S_1}(c_1),\,\dots,\,\mathrm{unflatten}_{S_m}(c_m)\big)
\]
for every coordinate; for $S' = S$ this says $f = S{:}D$. The question of
this section is whether $f$ is a layout $S'{:}D'$ for some refinement $S'$ of
$S$ and some stride $D'$.
\end{definition}

\begin{lemma}[Separability]\label{lem:separable}
Let $S = (n_1,\dots,n_m)$ be a flat shape and $f : \mathrm{Coord}(S) \to \Z$.
For $1 \le i \le m$ let $e_i$ be the coordinate with entry $i$ equal to $1$
and every other entry $0$, and define $g_i : [0,n_i) \to \Z$ by
$g_i(v) = f(v\,e_i)$. Then $f$ is a layout $S'{:}D'$ for some refinement $S'$
of $S$ if and only if
\begin{equation}\label{eq:sep}
f(c) = \sum_{i=1}^{m} g_i(c_i)
\quad\text{for every } c = (c_1,\dots,c_m) \in \mathrm{Coord}(S),
\end{equation}
and each $g_i$ is a layout of some flat shape of size $n_i$.
\end{lemma}
\begin{proof}
($\Rightarrow$) Let $f = S'{:}(D_1,\dots,D_m)$ with $S' = (S_1,\dots,S_m)$.
By Definition~\ref{def:layout},
\[
f(c) = \sum_{i=1}^{m} (S_i{:}D_i)\big(\mathrm{unflatten}_{S_i}(c_i)\big)
\quad\text{for every } c.
\]
At $c = v\,e_i$ every term but the $i$th is a layout evaluated at coordinate
$0$, hence $0$, so $g_i(v) = (S_i{:}D_i)(\mathrm{unflatten}_{S_i}(v))$: the
$i$th term is $g_i(c_i)$, which is \eqref{eq:sep}, and $g_i = S_i{:}D_i$.
($\Leftarrow$) If $g_i = S_i{:}D_i$ for each $i$, the same identity read
backwards gives $f = (S_1,\dots,S_m){:}(D_1,\dots,D_m)$.
\end{proof}

Checking \eqref{eq:sep} is one pass over $\mathrm{Coord}(S)$. What remains
is one coordinate at a time: a map $g : [0,n) \to \Z$ with $g(0) = 0$, and
the question whether there is a flat shape $T$ of size $n$ and a stride $D$
with $g = (T{:}D) \circ \mathrm{unflatten}_T$. To decide it we need this
composite as a formula in $v$.

Number the entries of $T$ from the right, $T = (m_k,\dots,m_1)$, so that
indices increase with the weights $w_j$ below. By Definition~\ref{def:layout},
$\mathrm{unflatten}_T(v)$ is the coordinate $(x_k,\dots,x_1)$ of $T$ with
\[
v = \sum_{j=1}^{k} x_j\, w_j, \qquad w_j = m_1 m_2 \cdots m_{j-1},
\]
the $w_j$ being the strides of $\mathrm{canonical}(T)$; so $w_1 = 1$,
$w_{j+1} = w_j m_j$, and $w_k m_k = n$. Solving for the components,
\[
x_j = \lfloor v / w_j \rfloor \bmod m_j,
\]
which we write $d_j(v)$. Therefore
\[
\big(T{:}(s_k,\dots,s_1)\big)\big(\mathrm{unflatten}_T(v)\big)
\;=\; \sum_{j=1}^{k} s_j\, d_j(v),
\]
and $g$ is the layout $T{:}(s_k,\dots,s_1)$ exactly when $g(v)$ equals this
sum for every $v \in [0,n)$.

Two remarks turn the search for $T$ into a search for the $w_j$. An entry
with $m_j = 1$ has $d_j = 0$ and contributes nothing, so we assume $T$ has
none; then $1 = w_1 < w_2 < \cdots < w_k < n$, each dividing the next and
$w_k$ dividing $n$, which we write $1 = w_1 \mid w_2 \mid \cdots \mid w_k
\mid n$ and call the \emph{chain} of $T$. Conversely, every such chain is
the chain of exactly one such $T$, given by $m_j = w_{j+1}/w_j$ with
$w_{k+1} = n$. Theorem~\ref{thm:decide} finds $T$ by finding its chain.

Each component $d_j(v) = \lfloor v/w_j \rfloor \bmod m_j$ is the remainder
of $\lfloor v/w_j \rfloor$ on division by $m_j$. The lemma rewrites
$\sum_j s_j\, d_j(v)$ with floors alone.

\begin{lemma}[Components as differences of floors]\label{lem:floor}
For positive integers $w \mid w'$ and every integer $v \ge 0$,
\[
\lfloor v/w \rfloor \bmod (w'/w) \;=\;
\lfloor v/w \rfloor - (w'/w)\lfloor v/w' \rfloor .
\]
Let $T$ be a flat shape of size $n$ with $k$ entries, with chain
$1 = w_1 \mid \cdots \mid w_k \mid n$, entry sizes $m_j$, and components
$d_j$ as above, and let integers $s_1,\dots,s_k$ and $a_1,\dots,a_k$ be
related by
\[
a_j = s_j - m_{j-1}\,s_{j-1} \quad (1 \le j \le k), \qquad m_0 s_0 = 0,
\]
a relation that determines either family from the other. Then for a map
$g : [0,n) \to \Z$,
\[
g(v) = \sum_{j=1}^{k} s_j\, d_j(v) \ \text{ for every } v \in [0,n)
\quad\Longleftrightarrow\quad
g(v) = \sum_{j=1}^{k} a_j \lfloor v/w_j \rfloor \ \text{ for every } v \in [0,n),
\]
and when these hold, $s_j = g(w_j)$.
\end{lemma}
\begin{proof}
Write $q = \lfloor v/w \rfloor$ and $r = w'/w$. The remainder of $q$ on
division by $r$ is $q - r\lfloor q/r \rfloor$, and $\lfloor q/r \rfloor =
\lfloor v/(wr) \rfloor = \lfloor v/w' \rfloor$ for positive integers $w, r$,
which is the identity. With $w = w_j$, $w' = w_{j+1}$, and $w'/w = m_j$ it
reads
\[
d_j(v) = \lfloor v/w_j \rfloor - m_j \lfloor v/w_{j+1} \rfloor,
\]
where for $j = k$, $\lfloor v/w_{k+1} \rfloor = \lfloor v/n \rfloor = 0$ on
$[0,n)$. Hence
\[
\sum_{j=1}^{k} s_j\, d_j(v)
= \sum_{j=1}^{k} s_j \lfloor v/w_j \rfloor
  - \sum_{j=1}^{k-1} s_j m_j \lfloor v/w_{j+1} \rfloor
= \sum_{j=1}^{k} \big(s_j - m_{j-1}s_{j-1}\big) \lfloor v/w_j \rfloor
= \sum_{j=1}^{k} a_j \lfloor v/w_j \rfloor ,
\]
which is the equivalence. The relation gives $s_j = a_j + m_{j-1}s_{j-1}$,
so either family determines the other. Finally,
$d_l(w_j) = \lfloor w_j/w_l \rfloor \bmod m_l$ is $0$ for $l > j$ because
$w_j < w_l$, is $1$ for $l = j$ because $m_j \ge 2$, and is $0$ for $l < j$
because $w_j/w_l$ is a multiple of $w_{l+1}/w_l = m_l$; so
$g(w_j) = \sum_l s_l\, d_l(w_j) = s_j$.
\end{proof}

With floors alone, first differences read the chain off $g$, because
$\lfloor t/w \rfloor - \lfloor (t-1)/w \rfloor$ is $1$ if $w \mid t$ and $0$
otherwise.

\begin{theorem}[Linear recognition]\label{thm:decide}
Let $n \ge 2$ (for $n = 1$, $g = 0$ is the layout $(1){:}0$), let
$g : [0,n) \to \Z$ with $g(0) = 0$, and let $h(t) = g(t) - g(t-1)$ for
$1 \le t < n$. Run the following scan. Keep a copy $\rho$ of $h$, an integer
$w_{\mathrm{last}}$ initially $1$, and a set $R$ initially empty. For
$t = 1, 2, \dots, n-1$ in order, if $\rho(t) \ne 0$: if $t \nmid n$ or
$w_{\mathrm{last}} \nmid t$, reject; otherwise put $t$ in $R$, set
$a_t = \rho(t)$ and $w_{\mathrm{last}} = t$, and replace $\rho(t')$ by
$\rho(t') - a_t$ for every multiple $t'$ of $t$ in $[t, n)$. If no $t$
rejects, accept.

Then $g$ is a layout of some flat shape of size $n$ if and only if the scan
accepts. If it accepts, $\{1\} \cup R$ is the chain of a flat shape $T$ of
size $n$; $g$ is the layout $T{:}(s_k,\dots,s_1)$ with $s_j = g(w_j)$; and
every flat shape of size $n$ without entries of size $1$ of which $g$ is a
layout has a chain containing $\{1\} \cup R$. We call this $T$ the
\emph{coarsest}. The scan takes $O(n)$ arithmetic operations.
\end{theorem}
\begin{proof}
By Lemma~\ref{lem:floor}, $g$ is a layout of the shape with chain $W$ if
and only if there are integers $a_w$, $w \in W$, with
$g(v) = \sum_{w \in W} a_w \lfloor v/w \rfloor$ for every $v \in [0,n)$.
Since $\lfloor t/w \rfloor - \lfloor (t-1)/w \rfloor = [\,w \mid t\,]$ and
$g(0) = 0 = \sum_{w} a_w \lfloor 0/w \rfloor$, this is equivalent to
\begin{equation}\label{eq:diff}
h(t) = \sum_{w \in W} a_w\, [\,w \mid t\,] \quad\text{for } 1 \le t < n.
\end{equation}

Suppose \eqref{eq:diff} holds for some chain $W$ and coefficients $a_w$. We
show by induction on $t$ that when the scan reaches $t$,
$R = \{\, w \in W : w < t,\ a_w \ne 0 \,\}$ with the recorded $a_w$ equal
to those of \eqref{eq:diff}, and $\rho(t) = a_t$ if $t \in W$ and
$\rho(t) = 0$ otherwise. At $t = 1$, $R$ is empty and $\rho(1) = h(1) =
a_1$ because $1 \in W$ is the only divisor of $1$. Given the claim at $t$:
$\rho(t) = h(t) - \sum_{w \in R,\, w \mid t} a_w = \sum_{w \in W,\, w \mid t}
a_w - \sum_{w \in W,\, w < t,\, w \mid t} a_w = [\,t \in W\,]\, a_t$. If
$\rho(t) \ne 0$ then $t \in W$ and $a_t \ne 0$; $t \mid n$ because every
element of $W$ divides $n$; and $w_{\mathrm{last}}$, which is the largest
element of $R$ or else $1$, lies in $W$ below $t$ and so divides $t$
because $W$ is a chain. The scan records $t$ with $a_t$, and the claim
holds at $t+1$; if $\rho(t) = 0$ nothing is recorded and the claim holds at
$t+1$ as well. So the scan accepts, and $R \subseteq W$; with $1 \in W$,
$\{1\} \cup R \subseteq W$.

Conversely suppose the scan accepts. At each $t$ either $\rho(t)$ was $0$
or $a_t = \rho(t)$ was subtracted from it, and later steps change $\rho$
only at or beyond their own $t$; so at the end $\rho \equiv 0$, that is,
$h(t) = \sum_{w \in R,\, w \mid t} a_w$ for $1 \le t < n$. Every element of
$R$ divides $n$, lies below $n$, and is a multiple of the element recorded
before it or of $1$, so $\{1\} \cup R$ is a chain; let $T$ be its shape,
and put $a_1 = 0$ if $1 \notin R$. Then \eqref{eq:diff} holds for
$W = \{1\} \cup R$, so $g$ is a layout of $T$ by Lemma~\ref{lem:floor}, with
strides $s_j = g(w_j)$.

Each element of $R$ is at least twice the one before, so the subtractions
touch fewer than $n + n/2 + n/4 + \cdots < 2n$ entries of $\rho$; the rest
of the scan is one pass over $t$.
\end{proof}

The methodology is similar to Appendix~A.2.1 of Axe~\cite{axe}.

Codegen therefore emits one of two expressions. If the map is a layout
$S'{:}D'$ for a refinement $S' = (S_1,\dots,S_m)$, it emits
$\sum_{i,j} s_{ij}\, d_{ij}(c_i)$, with $d_{ij}$ the components of
$\mathrm{unflatten}_{S_i}(c_i)$, for the coarsest $S'$; no such expression has
fewer terms. Otherwise it emits the composite as Definition~\ref{def:compose}
writes it: the value of the right operand, unflattened into the left operand's
domain shape, evaluated by the left operand. Section~\ref{sec:restrict}
extends this to restricted maps, which do not vanish at the origin.
